
\subsubsection{VAL Subcodes}
\label{sec:val-subcodes}

We identify six subcodes within the VAL family
across 119~GH statements and 4~KG statements.

% TODO: incorporate this into findings
% I think the use of precondition assertions are interesting. E.g., GH227 is a
% precondition assertion defined inside a function that creates a visualization.
% Writing a test for such a function does not make sense,
% yet we might make decisions downstream using this visualization.
% In such a case, using a precondition assert that explicitly fails
% when there is a fault in the data prevents error being propagated.

\paragraph{VAL-SHAPE}

Shape validation is the most prevalent form of explicit validation in the corpus.
Practitioners use these statements to verify
that arrays and tensors carry the dimensions that downstream operations require.
A shape mismatch in a PyTorch or NumPy pipeline
either raises a cryptic runtime error far from its origin,
or silently produces wrong results through unintended broadcasting.
Both failure modes are difficult to diagnose after the fact.

Shape validation concentrates at two points in the pipeline.
At the model construction stage,
practitioners verify that
each layer of a neural network produces the expected output shape
before the next layer is wired in.

\begin{lstlisting}[caption={}, label={lst:61}]
assert output['x5'].size() == torch.Size([BATCH_SIZE, 512, 7, 7])
\end{lstlisting}

For example, Listing~\ref{lst:61} checks the shape
of the fifth feature map from a VGG16 encoder
immediately before it is connected to a segmentation decoder.
A similar pattern appears across GAN encoders,
transformer attention blocks,
and GRU sequence models,
where practitioners assert the shape of intermediate tensors at each architectural boundary.

At the evaluation stage,
shape validation primarily guards against misaligned prediction and label arrays.
Before computing metrics such as RMSE or F1,
practitioners verify that predictions and ground truth have the same length.

\begin{lstlisting}[caption={}, label={lst:62}]
assert len(predictions_array) == len(label_batch_test.numpy())
\end{lstlisting}

Listing~\ref{lst:62} is not a trivial check.
A dropped remainder batch,
an off-by-one error in a data loader,
or a silent failure during sampling
can each produce a length mismatch that either crashes the metric computation
or, worse, produces a numerically plausible but incorrect score.

A third recurring theme is data integrity after joins, splits, and transformations.
Practitioners assert that a train-test split accounts for all rows,
that a join did not silently add or drop records,
and that a preprocessing step such as \lstinline{QuantileTransformer} preserves the number of samples.
These statements reflect an awareness
that pandas operations can silently alter the size of a dataframe
in ways that are not immediately visible in a notebook output.

\paragraph{VAL-EQ}

Practitioners check that a specific value, configuration parameter, or computed result
matches a known expected value.
The expected value is always statically known at the time of writing the assertion,
and is derived either from domain knowledge, architectural constraints, or a prior specification.

A prominent theme is the enforcement of architectural constraints
that have no runtime enforcement in the framework.
For example, sinusoidal positional embeddings
require an even embedding dimension
because the vector is split in half between sine and cosine components.
This is validated using Listing~\ref{lst:4}
which checks that the embedding dimensions are even.
If \lstinline{dim} is odd, downstream dimension expectations will not hold.

\begin{lstlisting}[caption={}, label={lst:4}]
assert dim % 2 == 0
\end{lstlisting}

Similarly, Wide ResNet filter counts must remain even at each resolution level,
and certain ResNet depth configurations
must satisfy a modular arithmetic relationship
to accommodate the bottleneck block structure.
These constraints are derived from the mathematical structure of the architecture
and are invisible to the framework at instantiation time.
Without the assertion,
a violated constraint produces a shape error
deep inside the forward pass whose connection to the misconfigured hyperparameter is not obvious.

A second theme is schema validation on loaded data.
Practitioners check that a CSV file has the expected number of classes,
% NOTE: e.g., common to use \lstinline{y = df[:, -1]} to get the labels from the last column
that a required column appears in the correct position,
that an index contains no duplicates,
and that a normalization step produced a zero-sum result.
These checks reflect the reality that data loading involves
external files whose schema cannot be verified statically.

\begin{lstlisting}[caption={}, label={lst:118}]
assert_frame_equal( X_train_prepro_ks.to_pandas(), X_train_prod_prepro_ks.to_pandas())
\end{lstlisting}

A third theme is the detection of training-serving skew.
Listing~\ref{lst:118} verifies
that the offline preprocessing pipeline and the production preprocessing pipeline
produce identical training data after feature selection.
A divergence between the two would cause the deployed model
to operate on a different feature distribution than the one it was trained on,
which would result in a silent failure with no crash signal.

\paragraph{VAL-BOOL}

Boolean validation captures a wide range of preconditions
that do not reduce to a shape check or an exact equality.
These statements guard pipeline steps against inputs
that are structurally valid but semantically incorrect.

The most common theme is data quality enforcement before a downstream step
that cannot recover from a violated precondition.
Practitioners assert the absence of missing values before feature selection,
the absence of class imbalance before oversampling logic executes,
and the satisfaction of domain constraints before domain-specific computations begin.

\begin{lstlisting}[caption={}, label={lst:54}]
assert not predictors.isnull().values.any()
\end{lstlisting}

For example, Listing~\ref{lst:54} is used in a cardiovascular disease mortality pipeline
to prevent missing values
from silently corrupting feature selection algorithms
or cause crashes in estimators that assume complete data.
The assertion ensures that an imputation or dropping step
was completed correctly upstream.

A second theme is the enforcement of geometric and physical invariants
derived from domain knowledge.
An object detection notebook asserts that both coordinates of a bounding box
satisfy the correct ordering before computing intersection over union.
A physics domain notebook asserts that model predictions
fall within the interval \([0, 1]\)
as required by the problem formulation.
A word2vec implementation asserts that the number of context words sampled
does not exceed the size of the context window.
These constraints cannot be inferred from the data alone
and require knowledge of the problem domain.

\begin{lstlisting}[caption={}, label={lst:22}]
assert all((layer.weight.grad is None
    for layer in dict(model.backbone.named_modules()).values()
    if hasattr(layer, 'weight')))
\end{lstlisting}

A third theme is transfer learning setup validation.
Listing~\ref{lst:22} asserts that all gradient tensors in a frozen backbone are \lstinline{None}.
A backbone that is not fully frozen continues to update during fine-tuning,
which silently degrades transfer learning performance without raising any error.

\paragraph{VAL-APPROX}

Approximate validation appears in contexts
where exact equality is either mathematically inappropriate
or unachievably strict due to floating-point arithmetic.
Two distinct use cases emerge from the data.

The first is unit testing of numerical implementations.
Practitioners validate custom loss functions,
data transformation utilities,
and metric implementations
by comparing their output to a known reference value within a specified tolerance.

\begin{lstlisting}[caption={}, label={lst:374}]
assert np.isclose(ndcg([3, 1, 2], [3, 1, 2], 3), 1)
\end{lstlisting}

For example, Listing~\ref{lst:374} validates
that a custom NDCG implementation returns a perfect score of 1.0
when predicted and true rankings are identical.

The second use case is reproducibility verification.
Practitioners who implement a notebook to reproduce the results of a published paper
use approximate assertions to confirm that their experimental setup matches the reported values.

\begin{lstlisting}[caption={}, label={lst:300}]
assert np.allclose(
    sp.stats.pearsonr(all_cfd_pred, all_truth)[0],
    0.409, atol=0.001)
\end{lstlisting}

Listing~\ref{lst:300} confirms that a Pearson correlation coefficient
matches the value reported in the source paper
within a tolerance of 0.001.
The tolerance is deliberately chosen
to allow for minor numerical differences across platforms and library versions
while still catching a substantively incorrect result.

\paragraph{VAL-EXIST}

Existence validation checks for membership in,
or presence within,
a dynamic runtime structure.
Two themes emerge.

The first is data leakage prevention.
In Listing~\ref{lst:84},
practitioners verify that training and validation sets are disjoint before model fitting begins.

\begin{lstlisting}[caption={}, label={lst:84}]
assert set(train_dataset.filenames).isdisjoint(
    set(valid_dataset.filenames))
\end{lstlisting}

Data leakage from validation into training inflates evaluation metrics
without any runtime error.
The only way to detect it programmatically
is to check for overlap explicitly,
as demonstrated in the listing.

The second theme is precondition checking on external resources.
Practitioners assert that a required file path exists
before attempting to load it,
and that a hyperparameter name is present in the parameter space of an estimator
before launching an expensive grid search.
In both cases the failure mode without the assertion
is either a cryptic I/O error
or a silent search over the wrong parameter space.

\paragraph{VAL-TYPE}

Practitioners use these statements to enforce
that objects carry the correct Python type or NumPy dtype
before passing them to operations that would fail silently
or with a misleading error
if given the wrong type.
Listing~\ref{lst:232} checks that model parameters are NumPy arrays
rather than Python lists,
since many NumPy operations accept both but produce different results.

\begin{lstlisting}[caption={}, label={lst:232}]
assert isinstance(mu, np.ndarray) and isinstance(sigma, np.ndarray),
    'mu/sigma must be numpy arrays.'
\end{lstlisting}

The low frequency of VAL-TYPE likely reflects the fact
that Python's dynamic typing and framework-level type coercion
often suppress type errors silently.
When a wrong type does not immediately crash the computation,
practitioners have no obvious signal that a check is needed.

\subsubsection{EXP Subcodes}
\label{sec:exp-subcodes}

We identify four subcodes within the EXP family
across 263~GH statements and 429~KG statements.
Unlike VAL statements, which are distributed relatively evenly between sources,
EXP statements are substantially more frequent in KG notebooks (62.0\% of all EXP statements).

\paragraph{EXP-INSPECT}

Inspection statements are the most frequent form of exploratory feedback in the corpus.
A practitioner places a bare variable name,
a plot call,
or a lightweight accessor expression as the last statement in a cell
to surface the current state of a data structure
without transforming it.
The output is consumed visually and immediately.

Inspection is used throughout the pipeline
but concentrates at the data loading and preparation stages,
where practitioners establish a running mental model of the dataset
before any transformation is applied.
A bare \lstinline{.head()} call after loading a CSV,
such as the one shown in Listing~\ref{lst:209},
confirms that column names parsed correctly,
that data types look plausible,
and that no obvious loading artefact is present.
The same lightweight pattern is used to monitor the effect of individual cleaning operations.
After imputing missing values,
a practitioner inspects the result to confirm that the fill worked.

\begin{lstlisting}[caption={}, label={lst:209}]
df_subjects.head()
\end{lstlisting}

At the evaluation stage,
inspection shifts from data structure verification to model output monitoring.
Practitioners surface intermediate predictions,
confusion matrices,
learned coefficients,
and cross-validation results
as bare last-cell expressions.

A recurring pattern during model training
is the periodic visual inspection of generated outputs or loss curves.
One notebook displays a GAN-generated image
at regular intervals during training
to assess whether the generator is producing realistic samples.
Another plots the training log-likelihood over epochs
to monitor convergence.
These statements serve as progress signals that are cheaper
and more immediate than a formal metric computation.

A distinct theme in GH notebooks
is the use of inspection to debug internal model state.
Practitioners surface the names of specific layers in a neural network,
the values of learned weight slices,
or the output of a forward hook
to understand what the model is doing internally.

\begin{lstlisting}[caption={}, label={lst:302}]
weights_conv1[:, :, 0, 0]
\end{lstlisting}

Listing~\ref{lst:302} slices a 4D convolutional weight tensor
to expose a single filter's learned spatial kernel,
which is then visualized to check what the filter has learned.
Inspection of this kind is not possible with a high-level metric,
as it requires looking directly at model internals.

\paragraph{EXP-COMPUTE}

Compute statements execute a transformation,
aggregation,
or model operation
and return a result that the practitioner examines.
The key distinction from inspection is that the result is not
already present in a variable,
rather it is produced on the fly by the statement itself.

The dominant use of EXP-COMPUTE in both GH and KG notebooks
is metric computation at the evaluation stage.
Practitioners compute accuracy, RMSE, F1, NDCG, AUC, and log-loss
as last-cell expressions
to obtain an immediate reading of model performance.

\begin{lstlisting}[caption={}, label={lst:339}]
accuracy_score(test_y, rf_pred_y)
\end{lstlisting}

These calls,
such as the one shown in Listing~\ref{lst:339}
are often repeated across multiple experimental runs,
each with a different model or hyperparameter setting,
to drive an iterative selection process.
One notebook computes cross-validated mean and standard deviation
of a Random Forest regressor's score
and compares it against a Logistic Regression
to decide which model to retain.

A second prominent theme is the use of EXP-COMPUTE
to check whether individual preprocessing steps worked as intended.
After applying a standard scaler,
a practitioner computes the mean and standard deviation of the transformed features
to verify that both are near their expected values (Listing~\ref{lst:293}).
After a \lstinline{dropna} operation,
a practitioner computes a null count to confirm that missing values were removed (Listing~\ref{lst:284}).

\begin{lstlisting}[caption={}, label={lst:293}]
(y_test.mean(), y_test.std())
\end{lstlisting}

\begin{lstlisting}[caption={}, label={lst:284}]
data.isnull().sum(axis=0)
\end{lstlisting}

These are informal postconditions.
They express an expectation about the output of the preceding operation,
but they rely on the practitioner's visual judgment
rather than a programmatic threshold.

A third theme,
distinctive to GH notebooks,
is the use of model fitting calls as pipeline integrity checks.
Practitioners run \lstinline{model.fit()} calls
and observe the training output
to confirm that the pipeline wiring is correct
before committing to a full training run.
One notebook runs a capsule network training call for 15 epochs
with a fraction of the data
to verify that loss decreases before launching the full experiment.

We code these statements as EXP-COMPUTE
because the fitting call is the statement under analysis
and the practitioner's intent is to observe whether the pipeline runs without error
rather than to record a final trained model.

In KG notebooks,
EXP-COMPUTE statements are also used extensively for exploratory feature analysis.
Practitioners compute grouped aggregations and correlation matrices
to understand whether a candidate feature
has a meaningful relationship with the target
before deciding whether to include it (Listing~\ref{lst:564}).

\begin{lstlisting}[caption={}, label={lst:564}]
train.drop('Type', axis=1).apply(lambda x: x.corr(train.Type))
\end{lstlisting}

\begin{lstlisting}[caption={}, label={lst:397}]
df_train.groupby(['Pclass', 'Sex'])['Survived'].value_counts(normalize=True)
\end{lstlisting}

Listing~\ref{lst:397} reveals a survival rate interaction between passenger class and sex.
The practitioner uses this insight
to motivate a subsequent feature engineering step.
This pattern of computing a grouped statistic
and acting on the result in the next cell,
is a core workflow unit in competition notebooks.

\begin{lstlisting}[caption={}, label={lst:480}]
test\_raw.shape[0] == test.shape[0]
\end{lstlisting}

A KG notebook uses Listing~\ref{lst:480}
as a last-cell boolean expression
to verify that no rows were lost during a transformation.
The expression is not wrapped in an \lstinline{assert},
so a False result would be visible in the cell output
but would not halt execution.
We code this statement as EXP-COMPUTE rather than VAL-SHAPE,
but the intent is identical to a shape assertion.

\paragraph{EXP-STRUCT}

Structural inspection statements
query the metadata of a data structure
rather than its values.
They surface shape, column names, dtypes, unique values, cardinality, or value counts.
EXP-STRUCT concentrates overwhelmingly at the preparation stage (78 of 126 statements),
where practitioners use these lightweight queries
to understand and verify the current state of the data
before applying the next transformation.

A pervasive theme is the use of structural inspection
as a lightweight checkpoint after each cleaning or encoding step.
After dropping columns,
a practitioner inspects \lstinline{.columns}
to confirm the drop worked.
After one-hot encoding,
a practitioner checks \lstinline{.shape}
to confirm that the column count increased as expected.
After undersampling the majority class,
a practitioner inspects the shape of the resulting arrays
to confirm that the class ratio changed (Listing~\ref{lst:593}).

\begin{lstlisting}[caption={}, label={lst:593}]
(X_train_under.shape, y_train_under.shape, X_test_under.shape)
\end{lstlisting}

These checkpoints serve the same purpose as postcondition assertions
but rely on visual confirmation
rather than programmatic enforcement.

A second theme is the use of \lstinline{.unique()} and \lstinline{.value\_counts()}
to audit the actual contents of a categorical column
before deciding on an encoding strategy.
Practitioners check whether unexpected values or label inconsistencies
are present before applying \lstinline{LabelEncoder} or \lstinline{OneHotEncoder}.

\begin{lstlisting}[caption={}, label={lst:528}]
combi['Item_Fat_Content'].value_counts()
\end{lstlisting}

The practitioner uses the output of Listing~\ref{lst:528}
to map all five variants of a categorical feature
to a binary encoding
before applying feature tools
for automated feature engineering.
Without this inspection,
the encoding library would treat the five strings as five distinct categories
and produce a semantically incorrect feature.

A third theme,
more common in GH notebooks,
is the use of structural queries to set downstream model parameters dynamically.
Practitioners inspect \lstinline{.vector\_size} of a word embedding
to determine the input layer dimension of a neural network,
or check \lstinline{X\_train.shape[2]}
to set the \lstinline{input\_shape} parameter of a 1D CNN.
These statements serve as interface checks
between the data preparation and model construction stages.

\paragraph{EXP-STATS}

Statistical summary statements
produce distributional views of a feature or target variable.
They include calls to \lstinline{.describe()},
\lstinline{.hist()},
\lstinline{.value\_counts()},
and distribution or box plots via seaborn and matplotlib.
EXP-STATS is proportionally more prevalent in KG notebooks (13.5\% of KG EXP)
than in GH notebooks (7.6\% of GH EXP),
which reflects Kaggle's culture of systematic exploratory data analysis before modelling.

The dominant theme is pre-modelling data auditing.
Practitioners use distribution plots and summary statistics
to detect outliers,
assess class imbalance,
and check whether a feature's distribution motivates a transformation.
A seaborn \lstinline{distplot} of insurance charges in Listing~\ref{lst:634}
reveals a right-skewed distribution,
which motivates a log transform
before fitting a linear regression model
whose residuals are assumed to be normally distributed.

\begin{lstlisting}[caption={}, label={lst:634}]
sns.distplot(data.charges)
\end{lstlisting}

A box plot of engine power values
in a used car dataset
drives a decision to remove extreme outliers.
A histogram of model prediction confidence scores
reveals whether the model is well-calibrated
or uncertain across its predictions.
In each case,
the statistical summary is not an end in itself
but a decision point that gates the next preprocessing or modelling step.

A second theme is post-processing verification.
After scaling features with \lstinline{StandardScaler},
practitioners call \lstinline{.describe()}
to confirm that the mean is near zero and the standard deviation is near one
(Listing~\ref{lst:580}).

\begin{lstlisting}[caption={}, label={lst:580}]
bikesData[columnsToScale].describe()
\end{lstlisting}

After imputing missing values for a column,
a practitioner plots the column's distribution
to confirm that the imputed values
do not distort the shape of the distribution
in a physiologically implausible direction.
These verification uses of EXP-STATS
occupy the same functional niche as EXP-STRUCT checkpoints,
by confirming that a preceding operation produced the expected result.
But they do so through distributional inspection
rather than structural metadata.

A third theme,
present primarily in GH notebooks,
is the use of statistical summaries to diagnose model behaviour.
A histogram of a model's confidence scores across a held-out test set
reveals whether predictions cluster near 0.5 (uncertain) or toward 0 and 1 (confident).
A KDE plot of SARIMA residuals checks whether the model's errors
approximate a normal distribution,
which is a formal assumption of the ARIMA model class.
These uses treat the statistical summary
as a model diagnostic tool rather than a data auditing tool.

